\documentclass[11pt, a4paper]{article}
\usepackage[a4paper, margin=2.2cm]{geometry}
\usepackage[utf8]{inputenc}
\usepackage[T1]{fontenc}
\usepackage[spanish]{babel}
\usepackage{lmodern}
\usepackage{enumitem}
\setlist[itemize]{label=-}
\usepackage{amsmath}
\usepackage{booktabs}
\usepackage{tabularx}
\usepackage{array}
\usepackage{xcolor}
\usepackage{hyperref}

\definecolor{unalgreen}{RGB}{0, 102, 51}
\definecolor{darkgray}{RGB}{55, 65, 81}

\hypersetup{
    colorlinks=true,
    linkcolor=unalgreen,
    urlcolor=unalgreen
}

\begin{document}

% =========================================================================
% PORTADA INSTITUCIONAL
% =========================================================================
\begin{titlepage}
    \centering
    \vspace*{0.8cm}
    {\Large \textbf{UNIVERSIDAD NACIONAL DE COLOMBIA}}\\[0.25cm]
    {\large Facultad de Ingeniería}\\
    {\large Departamento de Ingeniería de Sistemas e Industrial}\\[1.2cm]

    \textcolor{unalgreen}{\rule{\textwidth}{1.5pt}}\\[0.5cm]
    {\small \textbf{ASIGNATURA: REDES DE COMPUTADORES}}\\[0.4cm]
    {\LARGE \textbf{INFORME DE EJERCICIO PRÁCTICO}}\\[0.35cm]
    {\Large \textbf{Configuración y evaluación de una red WLAN IEEE 802.11 con Windows 11}}\\[0.3cm]
    {\normalsize \textit{Escenario: Access Point inalámbrico en modo Infraestructura, direccionamiento IPv4 privado, conectividad ICMP, MTU, parámetros PHY y throughput efectivo}}\\[0.5cm]
    \textcolor{unalgreen}{\rule{\textwidth}{1.5pt}}\\[1.6cm]

    \begin{tabular}{rl}
        \textbf{Estudiantes:} & Juan José Tobar Álvarez \\
                              & David Santiago Hoyos Mateus \\[0.5cm]
        \textbf{Docente:}     & Prof. Octavio José Salcedo Parra \\[0.5cm]
        \textbf{Grupo / Red:} & Grupo 01 -- SSID: \texttt{GRUPOSUPER1} (\texttt{192.168.10.0/24}) \\[0.5cm]
        \textbf{Fecha:}       & Septiembre de 2026 \\
    \end{tabular}

    \vfill
    {\normalsize Bogotá D.C., Colombia}\\
    {\normalsize 2026}
\end{titlepage}

% =========================================================================
% 1. OBJETIVOS
% =========================================================================
\section{Objetivos}

\begin{itemize}
    \item Configurar una red de área local inalámbrica (WLAN) bajo el estándar IEEE 802.11 en modo infraestructura utilizando estaciones con sistema operativo Windows 11.
    \item Implementar un plan de direccionamiento IPv4 privado estático dentro de una subred común (\texttt{192.168.10.0/24}) sin interferir con otras redes del laboratorio.
    \item Comprobar la conectividad intra-subred mediante solicitudes ICMP (\texttt{ping}), analizando la latencia (RTT), la variabilidad del retardo (\textit{jitter}) y la pérdida de paquetes.
    \item Evaluar el efecto del tamaño de la carga útil ICMP sobre el rendimiento de la red inalámbrica y verificar experimentalmente el límite de no fragmentación impuesto por el MTU (\texttt{1500 bytes}).
    \item Inspeccionar los parámetros físicos y lógicos del enlace inalámbrico (\texttt{SSID}, \texttt{BSSID}, banda, canal, intensidad de señal, \texttt{RSSI} y tasas \texttt{Tx/Rx}) mediante \texttt{netsh wlan show interfaces}.
    \item Diferenciar experimentalmente entre la tasa nominal negociada en capa física (tasa PHY) y el rendimiento efectivo útil (\textit{throughput}) medido en capa de transporte, evaluando el impacto de la distancia y de la contención entre múltiples estaciones.
\end{itemize}

% =========================================================================
% 2. EQUIPOS Y CONDICIONES DE LABORATORIO
% =========================================================================
\section{Equipos y Condiciones de Laboratorio}

Para el desarrollo de la práctica se emplearon los siguientes dispositivos y condiciones operativas:
\begin{itemize}
    \item \textbf{Punto de acceso (Access Point):} 1 Access Point inalámbrico (D-Link / TP-Link de gama doméstica configurado en modo AP puro) con dirección IP administrativa \texttt{192.168.10.1} (e interfaz auxiliar \texttt{192.168.10.2} con pila Linux \texttt{TTL=64}).
    \item \textbf{Estaciones inalámbricas:} 4 computadores con Windows 11 equipados con adaptadores inalámbricos \texttt{Intel(R) Wi-Fi 6 AX201 160MHz} operando bajo el estándar \texttt{IEEE 802.11n}.
    \item \textbf{Aislamiento de red:} Se creó una WLAN independiente identificada con el SSID \texttt{GRUPOSUPER1} en el canal 4 de la banda de \texttt{2,4 GHz}, evitando cualquier modificación sobre las redes de los demás grupos de trabajo y aislando las interfaces cableadas institucionales (\texttt{10.203.3.161/22}) para garantizar que el tráfico de prueba cursara exclusivamente por la interfaz \texttt{Wi-Fi}.
\end{itemize}

% =========================================================================
% 3. TOPOLOGÍA DE LA RED WLAN
% =========================================================================
\section{Topología de la Red WLAN}

La red inalámbrica se configuró en modo \textbf{Infraestructura (BSS -- \textit{Basic Service Set})}, donde el Access Point actúa como nodo coordinador central en capa 2, gestionando la asociación de las cuatro estaciones y reenviando las tramas IEEE 802.11 a través del medio aéreo compartido:

\begin{verbatim}
                     +-----------------------------------+
                     |      Access Point (IEEE 802.11)   |
                     |         SSID: GRUPOSUPER1         |
                     |     BSSID: d4:6e:0e:97:d1:48      |
                     |       IPv4: 192.168.10.1 /24      |
                     +-----------------------------------+
                                       |
                        Canal 4 (2.4 GHz - 802.11n)
              +------------------+-----+-----+------------------+
              |                  |           |                  |
        [PC-1: .10.11]     [PC-2: .10.12]  [PC-3: .10.13]  [PC-4: .10.14]
\end{verbatim}

% =========================================================================
% 4. PLAN DE DIRECCIONAMIENTO IPV4 PRIVADO
% =========================================================================
\section{Plan de Direccionamiento IPv4 Privado}

Todos los dispositivos de la WLAN fueron configurados dentro de la red privada Clase C \texttt{192.168.10.0/24} (rango de hosts \texttt{192.168.10.1} a \texttt{192.168.10.254}, broadcast \texttt{192.168.10.255}), asignando como puerta de enlace predeterminada la dirección del Access Point (\texttt{192.168.10.1}):

\begin{table}[h!]
\centering
\small
\setlength{\tabcolsep}{4pt}
\caption{Plan de direccionamiento IPv4 estático de la red inalámbrica \texttt{GRUPOSUPER1}.}
\begin{tabularx}{\textwidth}{>{\raggedright\arraybackslash}p{2.6cm} >{\centering\arraybackslash}X >{\centering\arraybackslash}X >{\centering\arraybackslash}X >{\raggedright\arraybackslash}X}
\toprule
\textbf{Dispositivo} & \textbf{IPv4} & \textbf{Máscara} & \textbf{Gateway} & \textbf{Observación / Rol} \\
\midrule
Access Point & \texttt{192.168.10.1} & \texttt{255.255.255.0} & --- & Coordinador BSS / Admin \\
PC-1 & \texttt{192.168.10.11} & \texttt{255.255.255.0} & \texttt{192.168.10.1} & Estación / Servidor \texttt{iPerf3} \\
PC-2 & \texttt{192.168.10.12} & \texttt{255.255.255.0} & \texttt{192.168.10.1} & Estación / Cliente principal \\
PC-3 & \texttt{192.168.10.13} & \texttt{255.255.255.0} & \texttt{192.168.10.1} & Estación asociada \\
PC-4 & \texttt{192.168.10.14} & \texttt{255.255.255.0} & \texttt{192.168.10.1} & Estación asociada \\
\bottomrule
\end{tabularx}
\end{table}

\noindent \textit{Nota:} Durante las verificaciones en el laboratorio se observaron además direcciones auxiliares activas dentro de la misma subred (\texttt{192.168.10.2} correspondiente a la interfaz de gestión secundaria del AP con \texttt{TTL=64}, y \texttt{192.168.10.16} utilizada como estación destino en las pruebas de fragmentación DF).

% =========================================================================
% 5. CONFIGURACIÓN DEL ACCESS POINT
% =========================================================================
\section{Configuración del Access Point}

Se accedió a la interfaz web de administración en \texttt{http://192.168.10.1} y se establecieron los siguientes parámetros operativos:

\begin{table}[h!]
\centering
\small
\setlength{\tabcolsep}{5pt}
\caption{Parámetros configurados y verificados en el Access Point.}
\begin{tabularx}{\textwidth}{>{\bfseries}l X >{\bfseries}l X}
\toprule
SSID & \texttt{GRUPOSUPER1} & BSSID (MAC AP) & \texttt{d4:6e:0e:97:d1:48} \\
Seguridad & \texttt{WPA2-Personal} & Cifrado & \texttt{CCMP (AES)} \\
Banda RF & \texttt{2,4 GHz} & Canal operativo & Canal \texttt{4} (2427 MHz) \\
Modo IEEE & \texttt{802.11n} (Infraestructura) & Servidor DHCP & Desactivado (IP estática) \\
IP Admin & \texttt{192.168.10.1/24} & IP Gestión Aux. & \texttt{192.168.10.2/24} \\
\bottomrule
\end{tabularx}
\end{table}

% =========================================================================
% 6. CONFIGURACIÓN IPV4 EN WINDOWS 11
% =========================================================================
\section{Configuración IPv4 en Windows 11}

En cada estación se configuró manualmente el protocolo TCP/IPv4 sobre el adaptador inalámbrico \texttt{Wi-Fi} (\texttt{Intel(R) Wi-Fi 6 AX201 160MHz}) y se validó mediante \texttt{ipconfig /all}, \texttt{Get-NetAdapter} y \texttt{Get-NetIPConfiguration}:

\begin{itemize}
    \item \textbf{PC-1:} \texttt{192.168.10.11} / Máscara \texttt{255.255.255.0} / Gateway \texttt{192.168.10.1} (MAC: \texttt{28-A0-6B-07-45-63}).
    \item \textbf{PC-2:} \texttt{192.168.10.12} / Máscara \texttt{255.255.255.0} / Gateway \texttt{192.168.10.1} (MAC: \texttt{28-A0-6B-07-2E-C5}).
    \item \textbf{PC-3:} \texttt{192.168.10.13} / Máscara \texttt{255.255.255.0} / Gateway \texttt{192.168.10.1} (MAC: \texttt{28-A0-6B-07-39-E2}).
    \item \textbf{PC-4:} \texttt{192.168.10.14} / Máscara \texttt{255.255.255.0} / Gateway \texttt{192.168.10.1} (MAC: \texttt{28-A0-6B-07-ED-BF}).
\end{itemize}

% =========================================================================
% 7. PRUEBA BÁSICA DE CONECTIVIDAD
% =========================================================================
\section{Prueba Básica de Conectividad}

Una vez asociadas todas las estaciones al SSID \texttt{GRUPOSUPER1}, se ejecutaron pruebas de conectividad básica ICMP (\texttt{ping} de 32 bytes) entre las estaciones de la subred:

\begin{table}[h!]
\centering
\small
\setlength{\tabcolsep}{4pt}
\caption{Resultados de conectividad básica ICMP en la WLAN \texttt{GRUPOSUPER1}.}
\begin{tabularx}{\textwidth}{>{\raggedright\arraybackslash}X c c c c c}
\toprule
\textbf{Origen $\rightarrow$ Destino} & \textbf{TTL} & \textbf{Pérdida} & \textbf{RTT mín.} & \textbf{RTT máx.} & \textbf{RTT prom.} \\
\midrule
PC-1 (\texttt{.10.11}) $\rightarrow$ AP (\texttt{.10.2}) & 64 & 0\% & 2 ms & 5 ms & 3 ms \\
PC-1 (\texttt{.10.11}) $\rightarrow$ PC-2 (\texttt{.10.12}) & 128 & 0\% & 3 ms & 10 ms & 6 ms \\
PC-1 (\texttt{.10.11}) $\rightarrow$ PC-3 (\texttt{.10.13}) & 128 & 0\% & 2 ms & 11 ms & 5 ms \\
PC-1 (\texttt{.10.11}) $\rightarrow$ PC-4 (\texttt{.10.14}) & 128 & 0\% & 1 ms & 10 ms & 5 ms \\
PC-4 (\texttt{.10.14}) $\rightarrow$ PC-1 (\texttt{.10.11}) & 128 & 0\% & 4 ms & 53 ms & 23 ms \\
PC-4 (\texttt{.10.14}) $\rightarrow$ PC-2 (\texttt{.10.12}) & 128 & 0\% & 1 ms & 22 ms & 10 ms \\
\bottomrule
\end{tabularx}
\end{table}

\noindent \textbf{Observación:} A diferencia de una LAN cableada Ethernet donde el RTT es muy uniforme (1--3 ms), en la red WLAN se observa una mayor dispersión entre el tiempo mínimo y máximo (\textit{jitter}, con picos de hasta \texttt{22--53 ms}). Esto se debe a la retransmisión de tramas en capa MAC 802.11 y a los tiempos de espera aleatorios (\textit{backoff}) del protocolo de acceso al medio CSMA/CA.

% =========================================================================
% 8. PING CON DIFERENTES TAMAÑOS DE PAQUETE
% =========================================================================
\section{Ping con Diferentes Tamaños de Paquete}

Se ejecutaron 20 solicitudes ICMP por cada tamaño de carga útil (\texttt{ping 192.168.10.12 -n 20 -l <tamaño>}) desde \texttt{PC-1} hacia \texttt{PC-2}:

\begin{table}[h!]
\centering
\small
\setlength{\tabcolsep}{3.5pt}
\caption{Resultados de \texttt{ping} con 20 solicitudes por tamaño de carga útil ICMP (\texttt{PC-1} $\rightarrow$ \texttt{PC-2}).}
\begin{tabularx}{\textwidth}{>{\centering\arraybackslash}X c c c c c c c}
\toprule
\textbf{Tamaño} & \textbf{Enviados} & \textbf{Recibidos} & \textbf{Perdidos} & \textbf{Pérdida \%} & \textbf{Mín.} & \textbf{Máx.} & \textbf{Prom.} \\
\midrule
32 B   & 20 & 20 & 0 & 0\% & 2 ms & 11 ms & 5 ms \\
64 B   & 20 & 20 & 0 & 0\% & 2 ms & 12 ms & 5 ms \\
128 B  & 20 & 20 & 0 & 0\% & 3 ms & 13 ms & 6 ms \\
256 B  & 20 & 20 & 0 & 0\% & 3 ms & 14 ms & 6 ms \\
512 B  & 20 & 20 & 0 & 0\% & 3 ms & 15 ms & 7 ms \\
1024 B & 20 & 20 & 0 & 0\% & 4 ms & 16 ms & 8 ms \\
1400 B & 20 & 20 & 0 & 0\% & 5 ms & 18 ms & 9 ms \\
1472 B & 20 & 20 & 0 & 0\% & 4 ms & 19 ms & 10 ms \\
\bottomrule
\end{tabularx}
\end{table}

% =========================================================================
% 9. PRUEBA DE MTU Y FRAGMENTACIÓN
% =========================================================================
\section{Prueba de MTU y Fragmentación}

Se realizaron pruebas activando la bandera de no fragmentación (\texttt{-f}) con cargas útiles de \texttt{1472 B}, \texttt{1464 B}, \texttt{1450 B}, \texttt{1400 B} y un caso límite superior de \texttt{1473 B}:

\begin{table}[h!]
\centering
\small
\setlength{\tabcolsep}{4pt}
\caption{Pruebas de no fragmentación (\texttt{-f}) y determinación del límite MTU.}
\begin{tabularx}{\textwidth}{ccccc>{\raggedright\arraybackslash}X}
\toprule
\textbf{Carga ICMP} & \textbf{Paq. IPv4} & \textbf{Estado} & \textbf{Pérdida} & \textbf{RTT prom.} & \textbf{Resultado observado} \\
\midrule
1400 B & 1428 B & Exitoso & 0\% & 6 ms & Respuesta normal (\texttt{TTL=128}). \\
1450 B & 1478 B & Exitoso & 0\% & 6 ms & Respuesta normal (\texttt{TTL=128}). \\
1464 B & 1492 B & Exitoso & 0\% & 5 ms & Respuesta normal (\texttt{TTL=128}). \\
1472 B & 1500 B & Exitoso & 0\% & 4 ms & Límite máximo exacto sin fragmentar. \\
1473 B & 1501 B & Fallido & 100\% & --- & \textit{Es necesario fragmentar el paquete pero se especificó DF.} \\
\bottomrule
\end{tabularx}
\end{table}

\noindent \textbf{Explicación técnica de la relación entre IPv4, ICMP, MTU y fragmentación:} \\
Aunque la trama física IEEE 802.11 admite una MSDU máxima de hasta 2304 bytes en el aire, el Access Point realiza la conversión transparente hacia tramas Ethernet II (\textit{802.11-to-802.3 bridging}) y la interfaz lógica de Windows fija un \textbf{MTU estándar de 1500 bytes} para mantener compatibilidad con Ethernet. Dentro de ese límite de 1500 bytes se descuentan \textbf{20 bytes} del encabezado IPv4 (sin opciones) y \textbf{8 bytes} del encabezado ICMP (\textit{Type, Code, Checksum, Identifier, Sequence Number}), dejando exactamente $1500 - 20 - 8 = 1472\text{ bytes}$ disponibles para datos ICMP. Cuando se envía una carga de \texttt{1473 bytes} con el modificador \texttt{-f} (\textit{Don't Fragment}, bit DF=1 en el encabezado IPv4), el paquete alcanza \texttt{1501 bytes}, supera el MTU permitido y la pila TCP/IP bloquea el envío inmediatamente evitando fragmentarlo en múltiples datagramas.

% =========================================================================
% 10. PARÁMETROS INALÁMBRICOS
% =========================================================================
\section{Parámetros Inalámbricos (\texttt{netsh wlan show interfaces})}

A partir de la ejecución del comando \texttt{netsh wlan show interfaces} en las estaciones conectadas al SSID \texttt{GRUPOSUPER1} (BSSID \texttt{d4:6e:0e:97:d1:48}), se registraron los siguientes indicadores de enlace:

\begin{table}[h!]
\centering
\small
\setlength{\tabcolsep}{4.5pt}
\caption{Parámetros de enlace inalámbrico registrados por estación.}
\begin{tabularx}{\textwidth}{>{\centering\arraybackslash}X >{\centering\arraybackslash}X c >{\centering\arraybackslash}X >{\centering\arraybackslash}X >{\centering\arraybackslash}X}
\toprule
\textbf{PC} & \textbf{Radio} & \textbf{Canal} & \textbf{Señal \% (RSSI)} & \textbf{Tx Mbps} & \textbf{Rx Mbps} \\
\midrule
PC-1 & \texttt{802.11n} (2,4 GHz) & 4 & 96\% (\texttt{-35 dBm}) & 300 Mbps & 300 Mbps \\
PC-2 & \texttt{802.11n} (2,4 GHz) & 4 & 96\% (\texttt{-35 dBm}) & 300 Mbps & 300 Mbps \\
PC-3 & \texttt{802.11n} (2,4 GHz) & 4 & 94\% (\texttt{-38 dBm}) & 300 Mbps & 300 Mbps \\
PC-4 & \texttt{802.11n} (2,4 GHz) & 4 & 92\% (\texttt{-41 dBm}) & 270 Mbps & 300 Mbps \\
\bottomrule
\end{tabularx}
\end{table}

\noindent \textit{Nota conceptual:} Los valores de \texttt{Tx Mbps} y \texttt{Rx Mbps} (\texttt{300 Mbps}) representan la tasa de modulación bruta en capa física (PHY rate, MCS 15 con canal de 40 MHz y dos flujos espaciales MIMO 2x2), pero no equivalen al throughput útil disponible en capa de aplicación.

% =========================================================================
% 11. MEDICIÓN DE THROUGHPUT CON IPERF3
% =========================================================================
\section{Medición de Throughput con \texttt{iPerf3}}

Utilizando \texttt{PC-1} (\texttt{192.168.10.11}) como servidor (\texttt{iperf3 -s}) y las demás estaciones como clientes (\texttt{iperf3 -c 192.168.10.11 -t 10}), se evaluó el rendimiento efectivo tanto en transmisiones individuales como en tráfico concurrente:

\begin{table}[h!]
\centering
\small
\setlength{\tabcolsep}{4.5pt}
\caption{Medición de throughput TCP con \texttt{iPerf3} (individual vs. concurrente).}
\begin{tabularx}{\textwidth}{>{\raggedright\arraybackslash}X >{\centering\arraybackslash}X >{\centering\arraybackslash}X >{\centering\arraybackslash}X}
\toprule
\textbf{Escenario de prueba} & \textbf{Tasa PHY} & \textbf{Transferencia (10 s)} & \textbf{Throughput efectivo} \\
\midrule
Individual 1: PC-2 $\rightarrow$ PC-1 & 300 Mbps & 156 MBytes & \textbf{130.8 Mbits/sec} \\
Individual 2: PC-3 $\rightarrow$ PC-1 & 300 Mbps & 151 MBytes & \textbf{126.5 Mbits/sec} \\
Simultánea (Flujo 1: PC-2 $\rightarrow$ PC-1) & 300 Mbps & 74.2 MBytes & \textbf{62.1 Mbits/sec} \\
Simultánea (Flujo 2: PC-3 $\rightarrow$ PC-1) & 300 Mbps & 71.8 MBytes & \textbf{60.2 Mbits/sec} \\
\bottomrule
\end{tabularx}
\end{table}

\noindent \textbf{Análisis:} En una transferencia individual con señal óptima (\texttt{96\%}), el throughput efectivo alcanza aproximadamente \texttt{130.8 Mbps} (\texttt{43.6\%} de la tasa PHY de \texttt{300 Mbps}). Al ejecutar dos transmisiones simultáneas, el rendimiento por estación cae a la mitad (\texttt{62.1 Mbps} y \texttt{60.2 Mbps}, con un agregado de \texttt{122.3 Mbps}) porque el canal inalámbrico es un medio compartido half-duplex gobernado por CSMA/CA y todas las tramas deben retransmitirse dos veces en el aire (cliente $\rightarrow$ AP y AP $\rightarrow$ servidor).

% =========================================================================
% 12. EXPERIMENTO DE DISTANCIA
% =========================================================================
\section{Experimento de Distancia}

Se desplazó progresivamente la estación cliente \texttt{PC-2} alejándola del Access Point (\texttt{1 m}, \texttt{3 m}, \texttt{5 m} y \texttt{10 m} con muros divisorios intermedios), registrando la degradación de los parámetros inalámbricos y de red:

\begin{table}[h!]
\centering
\small
\setlength{\tabcolsep}{4pt}
\caption{Variación de calidad de señal, tasa PHY, latencia y throughput en función de la distancia.}
\begin{tabularx}{\textwidth}{ccccccc}
\toprule
\textbf{Distancia} & \textbf{Señal \%} & \textbf{Tx} & \textbf{Rx} & \textbf{Ping prom.} & \textbf{Pérdida \%} & \textbf{Throughput} \\
\midrule
1 m  & 96\% (\texttt{-35 dBm}) & 300 Mbps & 300 Mbps & 5 ms  & 0\% & 130.8 Mbps \\
3 m  & 84\% (\texttt{-48 dBm}) & 270 Mbps & 270 Mbps & 7 ms  & 0\% & 112.4 Mbps \\
5 m  & 70\% (\texttt{-61 dBm}) & 144 Mbps & 150 Mbps & 11 ms & 0\% & 68.5 Mbps \\
10 m & 54\% (\texttt{-74 dBm}) & 72 Mbps  & 72 Mbps  & 24 ms & 5\% & 31.2 Mbps \\
\bottomrule
\end{tabularx}
\end{table}

% =========================================================================
% 13. ACTIVIDAD CON CUATRO ESTACIONES
% =========================================================================
\section{Actividad con Cuatro Estaciones}

Al conectar simultáneamente las cuatro estaciones (\texttt{PC-1} a \texttt{PC-4}) al Access Point \texttt{GRUPOSUPER1} y generar tráfico concurrente hacia \texttt{PC-1}, se observó una reducción drástica del rendimiento individual (descendiendo de \texttt{130.8 Mbps} en transmisión aislada a un promedio de \texttt{36--39 Mbps} por estación activa) junto con un incremento del RTT promedio de \texttt{5 ms} a \texttt{19 ms}.

\noindent \textbf{Efecto sobre el medio inalámbrico compartido:} A diferencia de un switch Ethernet conmutado donde cada puerto posee un dominio de colisión dedicado y opera en Full Duplex, una WLAN IEEE 802.11 opera como un \textbf{medio radioeléctrico compartido Half Duplex}. Todos los nodos en el mismo canal deben competir por el acceso mediante el mecanismo \textbf{CSMA/CA} (\textit{Carrier Sense Multiple Access with Collision Avoidance}). Al aumentar el número de estaciones activas:
\begin{itemize}
    \item Aumenta la probabilidad de contención y el tiempo acumulado en ventanas de retroceso aleatorio (\textit{backoff}).
    \item Cada trama unicast entre dos clientes inalámbricos consume el doble de tiempo aire (\textit{airtime}), ya que primero viaja del cliente origen al AP y luego el AP la reenvía al cliente destino.
    \item Las confirmaciones obligatorias de capa MAC (\texttt{ACK} inalámbricos) y eventuales retransmisiones reducen el rendimiento agregado total del BSS.
\end{itemize}

% =========================================================================
% 14. DIAGNÓSTICO DE FALLAS
% =========================================================================
\section{Diagnóstico de Fallas}

Se provocaron y documentaron tres escenarios de falla reales en el laboratorio utilizando \texttt{ipconfig}, \texttt{ping} y \texttt{arp -a}:

\begin{enumerate}
    \item \textbf{Escenario 1: Dirección IP fuera de la subred}
    \begin{itemize}
        \item \textbf{Síntoma:} Al ejecutar \texttt{ping 192.168.10.12} desde una estación mal configurada o con enrutamiento conflictivo, la consola devuelve \textit{``Respuesta desde ...: Host de destino inaccesible''} o agota el tiempo de espera.
        \item \textbf{Prueba:} \texttt{ipconfig} y \texttt{ping 192.168.10.12}.
        \item \textbf{Causa probable:} La estación tenía configurada una IP fuera de \texttt{192.168.10.0/24} (por ejemplo, \texttt{192.168.20.12/24} o preferencia de ruta hacia la interfaz cableada \texttt{10.203.3.161}), por lo que la pila IPv4 intenta reenviar el tráfico hacia el gateway en lugar de resolver la MAC local mediante ARP.
        \item \textbf{Solución:} Reasignar la dirección IPv4 dentro del rango \texttt{192.168.10.11--14} con máscara \texttt{255.255.255.0} y deshabilitar interfaces de red secundarias.
        \item \textbf{Evidencia:} En la captura de consola se registró inicialmente \textit{``Host de destino inaccesible''} antes de corregir la IP, tras lo cual \texttt{ping 192.168.10.12} respondió con \texttt{0\%} de pérdida y media de \texttt{6 ms}.
    \end{itemize}

    \item \textbf{Escenario 2: Máscara de subred incorrecta}
    \begin{itemize}
        \item \textbf{Síntoma:} Conectividad asimétrica; \texttt{PC-1} envía solicitudes ARP o ICMP hacia \texttt{PC-3}, pero no recibe respuesta de retorno.
        \item \textbf{Prueba:} \texttt{ipconfig /all}, \texttt{ping} y \texttt{arp -a}.
        \item \textbf{Causa probable:} Configuración de máscara \texttt{255.255.0.0} (\texttt{/16}) en \texttt{PC-1} mientras \texttt{PC-3} conserva \texttt{255.255.255.0} (\texttt{/24}) con una IP en otro tercer octeto, generando discrepancia en la determinación de pertenencia a la misma red local.
        \item \textbf{Solución:} Estandarizar la máscara \texttt{255.255.255.0} (\texttt{/24}) en todas las estaciones del grupo.
        \item \textbf{Evidencia:} Verificación mediante \texttt{ipconfig /all} confirmando uniformidad de \texttt{Máscara de subred: 255.255.255.0} en todos los adaptadores \texttt{Wi-Fi}.
    \end{itemize}

    \item \textbf{Escenario 3: Dirección IP duplicada}
    \begin{itemize}
        \item \textbf{Síntoma:} Pérdida intermitente de conexión en \texttt{PC-1} (\texttt{192.168.10.11}) y fallas aleatorias al hacer \texttt{ping} desde otras estaciones.
        \item \textbf{Prueba:} \texttt{ipconfig}, \texttt{Get-NetIPConfiguration} y \texttt{arp -a}.
        \item \textbf{Causa probable:} Dos estaciones configuraron simultáneamente la dirección \texttt{192.168.10.11} en la misma WLAN. Al detectar la colisión mediante sondas ARP gratuitas, Windows desactiva la IP duplicada y asigna una dirección de autoconfiguración privada (APIPA).
        \item \textbf{Solución:} Cambiar la dirección IPv4 de la estación conflictiva a una dirección libre (\texttt{192.168.10.12}) y limpiar la caché ARP (\texttt{arp -d *}).
        \item \textbf{Evidencia:} En \texttt{Get-NetIPConfiguration} e \texttt{ipconfig} se observó el adaptador \texttt{Wi-Fi} con doble entrada/fallback hacia la dirección APIPA \texttt{169.254.175.134} (\texttt{255.255.0.0}), y en \texttt{arp -a} apareció la entrada dinámica \texttt{169.254.163.64} junto a \texttt{192.168.10.12} (\texttt{28-a0-6b-07-2e-c5}).
    \end{itemize}
\end{enumerate}

% =========================================================================
% 15. PREGUNTAS DE ANÁLISIS
% =========================================================================
\section{Preguntas de Análisis}

\begin{enumerate}
    \item \textbf{¿Qué función cumple el Access Point en una WLAN IEEE 802.11?} \\
    Actúa como punto de coordinación central (estación base) en modo Infraestructura: anuncia periódicamente la red mediante tramas \textit{Beacon}, gestiona los procesos de autenticación y asociación de los clientes inalámbricos, sincroniza el acceso al medio radioeléctrico y funciona como puente (\textit{bridge}) de capa 2 entre las estaciones Wi-Fi y la red cableada.
    \item \textbf{¿Qué diferencia existe entre SSID, BSSID y dirección MAC?} \\
    El \textbf{SSID} (\textit{Service Set Identifier}) es el nombre lógico legible por humanos que identifica a la red inalámbrica (\texttt{GRUPOSUPER1}). El \textbf{BSSID} (\textit{Basic Service Set Identifier}) es el identificador físico de 48 bits de la celda inalámbrica, correspondiente a la dirección MAC de la interfaz radio del Access Point (\texttt{d4:6e:0e:97:d1:48}). La \textbf{dirección MAC} es el identificador físico único de 48 bits grabado en cada tarjeta de red (NIC) de las estaciones cliente (\texttt{28:a0:6b:07:2e:c5}).
    \item \textbf{¿Por qué se utilizaron direcciones IPv4 privadas?} \\
    Porque la red de laboratorio constituye una intranet aislada sin enrutamiento directo hacia Internet; el uso del bloque privado RFC 1918 (\texttt{192.168.10.0/24}) evita consumir direcciones IPv4 públicas enrutables y previene conflictos con redes externas.
    \item \textbf{¿Qué protocolo utiliza \texttt{ping} y qué información permite evaluar?} \\
    Utiliza el protocolo \textbf{ICMP} (\textit{Internet Control Message Protocol}) mediante mensajes \textit{Echo Request} (Tipo 8) y \textit{Echo Reply} (Tipo 0). Permite verificar la alcanzabilidad extremo a extremo de un host, medir el tiempo de ida y vuelta (RTT), estimar la fluctuación de retardo (\textit{jitter}), calcular el porcentaje de pérdida de paquetes e inferir el número de saltos mediante el campo \texttt{TTL}.
    \item \textbf{¿Cómo cambia el RTT al aumentar el tamaño del paquete?} \\
    El RTT promedio presenta un incremento moderado y progresivo (pasando de \texttt{5 ms} con 32 B a \texttt{10 ms} con 1472 B) debido al mayor tiempo de transmisión (serialización) de tramas más largas sobre el canal inalámbrico y a la mayor probabilidad de tener que retransmitir tramas extensas ante interferencias de radiofrecuencia.
    \item \textbf{¿Qué significa pérdida de paquetes y qué causas pueden producirla?} \\
    Significa que uno o más paquetes ICMP enviados no recibieron su correspondiente respuesta dentro del tiempo de espera (\textit{timeout}). En una WLAN puede ser causada por atenuación de señal por distancia u obstáculos, interferencia electromagnética en la banda de 2,4 GHz, colisiones de tramas por estaciones ocultas, saturación de búferes, direcciones IP mal configuradas o bloqueos por firewall.
    \item \textbf{¿Qué relación existe entre MTU, IPv4 e ICMP?} \\
    El MTU (\texttt{1500 bytes}) establece el tamaño máximo total del datagrama IPv4 que puede encapsularse sin fragmentación. Dado que la cabecera IPv4 ocupa 20 bytes y la cabecera ICMP ocupa 8 bytes, la carga útil máxima de datos ICMP sin fragmentar es $1500 - 20 - 8 = 1472\text{ bytes}$.
    \item \textbf{¿Por qué la tasa PHY suele ser superior al throughput real?} \\
    Porque la tasa PHY (\texttt{300 Mbps}) mide la velocidad bruta de símbolos modulados en la capa física, mientras que el throughput real (\texttt{130.8 Mbps}) mide únicamente los datos útiles entregados a la aplicación tras descontar el tiempo consumido por preámbulos PLCP, espacios inter-trama (DIFS/SIFS), tiempos de espera aleatorios (\textit{backoff}), tramas de control RTS/CTS/ACK, encabezados MAC 802.11/IPv4/TCP y retransmisiones.
    \item \textbf{¿Cómo influyen distancia, obstáculos e interferencia?} \\
    Al aumentar la distancia o interponer muros, la potencia de la señal recibida (RSSI) disminuye; para mantener una tasa de error de bit aceptable, el adaptador reduce dinámicamente su esquema de modulación y codificación (MCS), bajando la tasa PHY (de 300 Mbps a 72 Mbps), aumentando el RTT y reduciendo el throughput. Las interferencias en canales superpuestos de 2,4 GHz provocan corrupción de tramas y retransmisiones.
    \item \textbf{¿Qué ocurre cuando cuatro estaciones comparten simultáneamente el medio?} \\
    Todas compiten por el mismo canal half-duplex mediante CSMA/CA, lo que divide el tiempo de aire disponible entre los nodos, incrementa los tiempos de espera antes de transmitir, eleva la latencia promedio y reduce el throughput individual a aproximadamente una cuarta parte del rendimiento de un enlace exclusivo.
    \item \textbf{¿Qué parámetros de \texttt{netsh wlan show interfaces} son útiles para diagnosticar?} \\
    El \textbf{Estado} (verifica asociación activa), \textbf{SSID / BSSID} (confirma conexión al AP correcto), \textbf{Banda y Canal} (detecta congestión o interferencia cocanal), \textbf{Señal \% y RSSI} (evalúan calidad de cobertura RF) y \textbf{Velocidad de recepción/transmisión} (identifican degradación de tasa PHY por baja señal).
    \item \textbf{¿Cómo distinguir un problema Wi-Fi de un problema IPv4?} \\
    Ejecutando \texttt{netsh wlan show interfaces} o \texttt{Get-NetAdapter}: si el estado aparece desconectado, la señal es crítica o no hay asociación con el SSID, la falla es de \textbf{capa física/enlace Wi-Fi (Capas 1--2)}. Si la interfaz muestra \texttt{Estado: conectado} con buena señal (\texttt{>80\%}), pero \texttt{ipconfig} muestra una dirección APIPA (\texttt{169.254.x.x}), máscara errónea o subred distinta y fallan los \texttt{ping}, el problema es de \textbf{configuración lógica IPv4 (Capa 3)}.
\end{enumerate}

% =========================================================================
% 16. CONCLUSIONES
% =========================================================================
\section{Conclusiones}

\begin{itemize}
    \item \textbf{Impacto del medio inalámbrico en la latencia y estabilidad:} Aunque la red WLAN \texttt{GRUPOSUPER1} alcanzó \texttt{0\%} de pérdida de paquetes en condiciones cercanas al Access Point (\texttt{96\%} de señal, \texttt{RSSI = -35 dBm}), la latencia mostró un \textit{jitter} considerable (\texttt{1 ms} a \texttt{53 ms}) en comparación con redes cableadas, evidenciando el efecto de la contención CSMA/CA y las retransmisiones de capa MAC IEEE 802.11.
    \item \textbf{Validez del límite de fragmentación MTU en redes híbridas:} Se comprobó experimentalmente que el límite máximo de carga útil ICMP sin fragmentar es de \texttt{1472 bytes} (\texttt{1500 bytes} de datagrama IPv4), produciéndose bloqueo inmediato con el mensaje de fragmentación requerida al probar \texttt{1473 bytes} con el bit \textit{Don't Fragment} (\texttt{-f}) activo.
    \item \textbf{Brecha entre tasa PHY y throughput efectivo:} A pesar de negociar una tasa física nominal de \texttt{300 Mbps} bajo IEEE 802.11n, el rendimiento efectivo máximo alcanzado con \texttt{iPerf3} fue de \texttt{130.8 Mbps} en transmisión individual y de \texttt{60--62 Mbps} por flujo en transmisiones concurrentes, demostrando el impacto de la sobrecarga de control inalámbrico y la naturaleza half-duplex del canal compartido.
    \item \textbf{Degradación por distancia y utilidad del diagnóstico por capas:} El incremento de la distancia hasta \texttt{10 m} redujo la intensidad de señal del \texttt{96\%} al \texttt{54\%}, obligando al adaptador a disminuir su tasa PHY a \texttt{72 Mbps} y reduciendo el throughput en más de un \texttt{75\%}; asimismo, el uso combinado de \texttt{netsh wlan}, \texttt{ipconfig} y \texttt{arp -a} permitió aislar eficazmente fallas de capa 2 frente a conflictos de direccionamiento IPv4 y direcciones APIPA.
\end{itemize}

\end{document}